\chapter{Conclusion and Future Work}
\label{ch:conclusion}

\section{Summary of Findings}
\label{sec:summary}

I did a comparison of five machine learning classifiers that predict diseases based on binary symptom vectors: Naive Bayes, Logistic Regression, SVM, Random Forest, and XGBoost. I ran each of them on two Kaggle datasets: one with approximately 5,000 records and one with approximately 246,000 records with the same preprocessing, cross-validation and evaluation processes. The aim was to give a controlled comparison to answer the question: what classifier is most effective to predict disease based on symptoms, and is the answer to this question sensitive to the size of the dataset or skewed class distributions?

The findings indicated that there were distinct performance variations, yet no model performed optimally across all contexts. In Dataset 1, it can be observed that the scores of Logistic Regression and SVM were perfect, and Naive Bayes and Random Forest were also good. Naive Bayes performed the highest accuracy on Dataset~2, and the Logistic Regression and XGBoost also performed well. In general, the optimal model is determined by the size of the dataset, the imbalance of classes, and the accuracy-training-interpretability trade-off.

Logistic Regression was still competitive, as a linear model. It was not as strong on every measure as the best models, yet remains an option when interpretability is valued more than using the final few percentage points of accuracy.

SVM was like Logistic Regression in predictive measures but it was much slower to train especially on the large dataset. One-vs-one multi-class strategy is not that scaled well with the number of classes hence not practical in this application.

Naive Bayes was a powerful baseline and it was particularly competitive on Dataset~2 where it recorded the highest accuracy. It continues to train very fast, thus making it still useful when the simplicity and runtime are more important than the model complexity.

The additional information shifted the image, although not in the same direction. There were models which took advantage of the increased dataset and others failed to become the best overall. This indicates that the actual data should form the basis of model choice when predicting diseases using symptoms, as opposed to making a general assumption that more sophisticated procedures cannot lose.

Error analysis showed that majority of the misclassifications occurred between diseases with similar symptoms. E.g. the various respiratory conditions were often mixed in that they share common symptoms such as cough, fever and chest pain. This is medically explainable and not necessarily an issue to the healthcare booking use case as patients having similar symptom patterns usually require the same kind of specialist anyway.

\section{Contributions}
\label{sec:contributions}

This thesis makes three contributions:

\begin{enumerate}
    \item A comparative evaluation of five classifiers to the same symptom based datasets with the same preprocessing and evaluation setup. The majority of current research experiments one or two classifiers in different settings and this means that comparison between papers cannot be reliable.
    \item A comparison on two datasets, of vastly different size, of the relationship between the performance of classifiers and the availability of data. This is what I could not find in the literature.
    \item An explanation of the steps needed to implement a symptom-based prediction model in a healthcare booking system, both in terms of regulatory obligations pursuant to the EU AI Act and the GDPR.
\end{enumerate}

\section{Recommendations}
\label{sec:recommendations}

According to my discovery, I suggest the following:

\begin{itemize}
    \item Logistic Regression is used when interpretability is important, since the coefficients are simple to interpret and describe.
    \item Use Naive Bayes as a robust and very fast baseline, particularly when there is a constraint on training time or when a simple reference model is required.
    \item XGBoost and Random Forest are helpful tools, and one should not expect them to always be more effective than simpler models on any given dataset.
    \item Instead of SMOTE, use class weighting to deal with the imbalanced classes when features are binary.
    \item Should not be deployed without external validation on actual clinical data, confidence thresholds, and human supervision.
\end{itemize}

\section{Future Work}
\label{sec:future_work}

The thesis will offer a basis on which disease can be predicted through symptoms, yet there are a number of significant directions that where the research can be directed towards in future.

\subsection{Clinical Validation}

The next critical thing is to test these classifiers on actual clinical data. The Kaggle datasets that I have used in this case are helpful in the controlled comparison, but they do not reflect the chaotic nature of the real patient encounters. The actual patients may lose the symptoms, misunderstand questions or give different reports than anticipated. It is quite possible that the disease distribution in an actual clinic is quite different than that in the training data.

It would be possible to collaborate with a hospital or clinic to acquire anonymized patient records, which would enable a much more realistic assessment. This would be based on the gathering of symptom data on real patient experiences (self-reported via a Web form or documented by the nurses at the intake) and the confirmed diagnosis by the doctor. This real-world data could then be trained on by the classifiers and tested on held-out cases.

A study like this would require an institutional review board to approve such a study and would be required to abide by GDPR and other regulations on data protection. It would also have to solve the problem that confirmed diagnoses may not be immediately available --- not all conditions can be diagnosed by lab tests or imaging, which may take days or weeks.

\subsection{Incorporating Additional Features}

The datasets that I employed only have binary symptom data. It is possible that adding patient demographics (age, sex, medical history) as other features would enhance the accuracy of the prediction. Of special importance is age since the prevalence of diseases differs significantly with age. The clinical picture of a rash in a child is not the same as that of the rash in an elderly.

History of medicine is good as well. A patient who has a history of heart disease and presents with chest pain should be diverted as opposed to a healthy 25-year-old reporting the same. Nonetheless, the addition of medical history presents privacy risks and makes the system more complex.

Symptom severity and duration may also be added as another way. The system would not need to binary encode (absent or present), but the patients would be asked to grade the severity on a scale (mild, moderate, severe) and specify the duration that they have had each symptom. It would give the classifier more information and would complicate and lengthen the questionnaire.

\subsection{Multi-Label Classification}

My datasets and methods are based on the assumption that a patient has no more than one disease. In practice, patients are frequently diagnosed with more than one condition (comorbidities). One may either be diabetic and infected with a skin infection, or asthmatic and allergic to seasons. Multi-label classification would enable the model to classify more than one disease of a patient.

This would involve training data in which every patient is classified with all of his or her conditions, not only the primary one. It would also involve changes in the evaluation measures and specialist-mapping logic. Should the model forecast a dermatological condition and a respiratory condition, which specialist should the patient make an appointment with? The system may have to rank according to severity or recommend appointing with several specialists.

\subsection{Specialist Mapping Layer}

I did not apply the disease-to-specialist mapping and concentrated on disease prediction. The creation of this layer would entail medical knowledge of what diseases are referable to which specialists. Certain mappings are self-evident (eczema dermatologist), whereas others are more complicated. A diabetic patient may require an endocrinologist whereas a diabetic patient may require a podiatrist in case of diabetes related foot issues.

The mapping could also rely on the condition severity. A GP may attend to a mild asthma case, whereas a pulmonologist is necessary in a severe case. The addition of severity assessment to the specialist recommendation logic is a promising trend to be pursued in the future.

\subsection{Reduced-Feature Models}

The analysis of feature importance (Section~\ref{sec:feature_importance}) demonstrated that there are symptoms that are far more informative than others. It may be possible to test a reduced-feature model based on only the top 15-20 most important symptoms to determine whether a simpler questionnaire can perform as well. It would make the system much more practical to the patients in case it does.

A 132 yes/no questionnaire is cumbersome and may result in low quality of the data since patients lose concentration. A reduced questionnaire containing 15-20 well selected symptoms would be easy to fill in and may even enhance the quality of data by minimizing the probability of patients selecting the questions by mistake or by clicking the questions randomly.

The issue is which symptoms should be retained. The rankings of the feature importance are global (average of all diseases), and some of the symptoms may be important in particular diseases, though they are not important in all diseases. It would require a careful analysis to identify a subset of the symptoms that would be well covered in all disease categories.

\subsection{Ensemble Methods and Model Stacking}

Another interesting direction is to combine several classifiers instead of choosing one. An example would be XGBoost as the main model, but when XGBoost and Logistic Regression disagree with one another on a prediction, this may indicate uncertainty, and the patient may be referred to a GP rather than a specialist.

More advanced ensemble techniques such as stacking (training a meta-model that takes the predictions of several base models) may theoretically be even more effective. This, however, complicates the system and makes it more difficult to interpret and maintain.

\subsection{Active Learning and Adaptive Questionnaires}

An adaptive questionnaire may provide follow-up questions based on the responses given earlier instead of questioning each patient on all 132 questions to assess their symptoms. As an example, when a patient presents with a rash, the system may ask questions related to the skin issue (itching, location, duration) but not irrelevant symptoms.

This would reduce the length of the questionnaire and make it more individualized, however, it would complicate the representation of the features and would necessitate an alternative form of modeling, which may be a decision tree or reinforcement learning.

\subsection{Deployment and Monitoring}

Lastly, the construction and implementation of a full system would entail all of the pragmatic aspects covered in Section~\ref{sec:deployment}: the need to integrate with appointment scheduling systems, confidence thresholds, edge cases, and performance monitoring. A pilot implementation in a small clinic or a university health center would be informative as to what works and what does not work in practice.